% =====================================================================
%  Draft section text for the DORI poster, as it stood before the blocks
%  were emptied. Paste each piece into the matching block.
%
%  TWO NUMBERS ARE STALE. This draft says a dead time of 18 us and an energy
%  range of 20-50 keV. Your table now says 17.8 +/- 0.4 us and 20 to 60 keV.
%  Fix them wherever they appear below before using it.
% =====================================================================


% --------------------------------------------------------- INTRODUCTION
Interventional radiology is the medical specialty subject to the highest
occupational exposure to ionising radiation, since it relies on
fluoroscopy-guided techniques [1]. Occupational monitoring currently rests on
passive dosimeters, which report dose only \emph{a posteriori}. The ability to
anticipate and to act upon unnecessary exposure is therefore limited [2].

This work proposes DORI, a personal and active \textbf{DO}simetry system for
\textbf{R}adiology of \textbf{I}ntervention. DORI targets real-time monitoring
of scattered radiation through a low-cost unit that can be replicated for
monitoring at chest, neck, hand and eye level.


% -------------------------------------------------------------- METHODS
Each unit integrates a plastic scintillator, chosen for its soft-tissue
equivalent properties, coupled to a silicon photomultiplier. Custom readout
electronics (charge-sensitive amplifier, pole-zero cancellation, shaper and
peak-and-hold) discriminate and digitise the signal amplitude. Wireless
communication is handled by a battery-powered microcontroller, keeping the
assembly compact and unobtrusive to the operator. Dose-rate data are received
by a Raspberry Pi and streamed in real time to an external visual interface.

\textbf{Energy calibration.} Bremsstrahlung endpoints (Duane--Hunt) acquired
between 20 and 50\,kVp, combined with the $^{241}$Am photopeak at 59.5\,keV in
a weighted least-squares fit. The $^{57}$Co Compton edge was reserved as an
independent validation point.

\textbf{Dose calibration.} ISO 4037 narrow-spectrum qualities N-25, N-30 and
N-40, with $\dot{H}_p(10)$ reconstructed from the measured spectrum through a
bin-weighted least-squares fit.


% -------------------------------------------------------------- RESULTS
The readout electronics were characterised and present a dead time of
18\,$\mu$s. DORI is calibrated over the energy range relevant to this
application (20--50\,keV), where it shows linear behaviour:
\[
  \mathrm{channel} = a\,E + b
\]
\vspace{-8mm}
\[
  a = (51\pm2)\times10^{-2}~\mathrm{channel\,keV^{-1}},\quad
  b = 21\pm1,\quad R^2 = 0.989 .
\]


% ----------------------------------------------------------- DISCUSSION
The response is linear over the energy range relevant to fluoroscopy-guided
procedures, so a single pair of calibration parameters describes the whole
working range. The dead time of 18\,$\mu$s is short compared with the count
rates expected from scattered radiation at the operator position.

The largest deviation occurs at N-25, the quality with the lowest air kerma
rate and therefore the largest measurement uncertainty.


% ---------------------------------------------------------- CONCLUSIONS
\begin{dorilist}
  \item Active, wireless and low-cost personal dosimetry is feasible with a
        plastic scintillator read out by a silicon photomultiplier.
  \item The prototype is calibrated in energy over 20--50\,keV, with linear
        response and $R^2 = 0.989$.
  \item A dead time of 18\,$\mu$s is compatible with the scattered fields
        found in interventional procedures.
  \item $\dot{H}_p(10)$ is reconstructed directly from the measured spectrum,
        without an auxiliary reference instrument.
\end{dorilist}


% ----------------------------------------------------------- REFERENCES
[1] B.\,A. Schueler and K.\,A. Fetterly, ``Eye protection in interventional
procedures,'' \emph{British Journal of Radiology}, vol.~94, p.~20210436, 2021.

[2] K. Kromp\ss{}a\ss{} \emph{et al.}, ``Real-time dosimetry in interventional
radiology: comparing the occupational radiation exposure in fluoroscopy guided
lower extremity and abdominal procedures,'' \emph{European Radiology}, 2025.
